\section{The FreshState Pipeline}
\label{sec:method}

\subsection{Problem Formulation}

Let $\URL$ be a set of URLs. For each URL $u \in \URL$ and timestamp $t$, let
$\snap{u}{t}$ denote the HTML content of $u$ at time $t$, and let $\val{u, t}$
denote the answer-bearing value extracted from $\snap{u}{t}$ by a domain-specific
extractor $f_d$:
\[
  \val{u, t} = f_d\!\left(\snap{u}{t}\right) \;\in\; \mathcal{V} \cup \{\bot\},
\]
where $\mathcal{V}$ is the value space (e.g., prices, availability labels) and
$\bot$ indicates extraction failure.

\begin{definition}[Change event]
  A \emph{change event} at URL $u$ between times $t_i < t_j$ is a triple
  $(u,\, t_i,\, t_j)$ such that $\val{u, t_i} \neq \val{u, t_j}$ and
  $\val{u, t_i} \neq \bot$ and $\val{u, t_j} \neq \bot$.
\end{definition}

\begin{definition}[Staleness label]
  Given a search result snippet $q$ cached at time $\tbase$ and served at time
  $\tcheck > \tbase$, the staleness label is $\stale(q) = 1$ if there exists a
  change event $(u, t_i, t_j)$ with $t_i \leq \tbase < t_j \leq \tcheck$, and
  $\stale(q) = 0$ otherwise.
\end{definition}

FreshState's goal is to produce a dataset of $(q,\, \stale(q))$ pairs, where
staleness labels are derived from automatically collected change events rather than
human annotation.

\subsection{Candidate Collection}

Before monitoring can begin, FreshState collects a set of \emph{candidate URLs}
$\URL^+$ that are likely to exhibit changes within the monitoring window. Two
strategies are supported.

\paragraph{Strategy A: Retrospective (Wayback CDX).}
Given a URL pattern (e.g., \texttt{www.bhphotovideo.com/c/product/*}), the
Wayback Machine CDX API is queried for all archived snapshots matching the pattern.
URLs that appear in multiple snapshots separated by at least $\lag_{\min}$ days are
retained as candidates, since their archive history implies they are stable enough
to persist through the monitoring period. This strategy requires no waiting---all
data is already in the archive---but depends on Wayback crawl coverage, which varies
by domain and is occasionally unavailable.

\paragraph{Strategy B: Prospective Monitor.}
Candidates are collected from live sources known to contain high-churn pages:
Craigslist housing listings (scraped via HTML) and GitHub release pages (identified
via the GitHub Search API for actively maintained repositories). A rolling collection
strategy adds new candidates daily while continuing to monitor all previously
baselined URLs. After collection, a \emph{baseline fetch} records $\snap{u}{\tbase}$
and $\val{u, \tbase}$ for every candidate. A daily scheduled job then fetches each
URL again, computes $\val{u, \tcheck}$, and records a change event if the value has
shifted. \Cref{alg:monitor} summarizes the per-URL monitoring step.

\begin{algorithm}[t]
\caption{Per-URL monitoring step (Strategy B)}
\label{alg:monitor}
\begin{algorithmic}[1]
\REQUIRE URL $u$, baseline value $v_0 = \val{u, \tbase}$, domain $d$, state file $S$
\STATE Fetch live HTML: $h \leftarrow \texttt{fetch\_live}(u)$
\IF{HTTP status $= 410$}
  \STATE Record expiration event $(u,\, \tbase,\, \tcheck)$ in $S$; \textbf{continue}
\ENDIF
\STATE Extract value: $(v, \text{span}, \conf) \leftarrow f_d(h)$
\IF{$\conf < \conf_{\min}$}
  \STATE Skip (low-confidence extraction)
\ELSIF{$v \neq v_0$}
  \STATE Append change event to seeds file; update $S[u] \leftarrow v$
\ENDIF
\end{algorithmic}
\end{algorithm}

\subsection{Domain-Specific Extractors}

FreshState provides a library of extractors, each implementing the interface
$f_d : \text{HTML} \to (\mathcal{V} \cup \{\bot\},\; \text{span},\; \conf)$.
The returned triple consists of a normalized value, the surrounding text span
(for interpretability and snippet alignment), and a confidence score $\conf \in [0, 1]$.

\paragraph{Price extractor.}
Targets rental and product pages. Uses a cascade of CSS selectors
(\texttt{[data-testid='price']}, \texttt{.listing-price}, \texttt{[itemprop='price']},
etc.) before falling back to a regex over raw page text:
\[
  \texttt{PRICE\_RE} = \texttt{\$\textbackslash s?[\textbackslash d,]+(?:\textbackslash
  .\textbackslash d\{2\})?(?:/mo)?}
\]
Selector matches yield $\conf = 0.9$; regex fallbacks yield $\conf = 0.6$.

\paragraph{Availability extractor.}
Matches patterns such as ``available now'', ``no longer available'', ``rented'',
and ``move-in date'' against page text, returning a categorical label from
$\mathcal{V}_{\text{avail}} = \{\texttt{available\_now},\, \texttt{available\_date},\,
\texttt{unavailable},\, \texttt{off\_market}\}$.

\paragraph{Version extractor.}
Targets software and product-spec pages. Matches version strings of the form
\texttt{v?X.Y.Z} using structured selectors and full-text regex.

\paragraph{Snippet builder.}
Independently of value extraction, FreshState builds a search-snippet representation
of each page by cascading through \texttt{<meta name="description">},
\texttt{<meta property="og:description">}, and the text surrounding the
answer-bearing span. This snippet is the unit that receives a staleness label,
mirroring the snippet a search engine would surface.

\subsection{Change Event Format}

Each change event is stored as a JSON line in a \emph{seeds file}:
\begin{verbatim}
{
  "url":        "https://sfbay.craigslist.org/...",
  "domain":     "apartment",
  "baseline_t": "2026-03-28T00:00:00Z",
  "check_t":    "2026-03-29T18:00:00Z",
  "v_old":      "$2400/mo",
  "v_new":      "$2350/mo",
  "span":       "...$2350/mo... available now ...",
  "conf":       0.9,
  "change_type":"price_change"
}
\end{verbatim}
Expiration events (HTTP 410) are stored with \texttt{v\_new = null} and
\texttt{change\_type = "expiration"}.
